\documentclass[11pt]{article}
\usepackage[utf8]{inputenc}
\usepackage[T1]{fontenc}
\usepackage{natbib}
\usepackage{times}
\usepackage{latexsym}
\usepackage{microtype}
\usepackage{hyperref}
\bibliographystyle{plainnat}

\title{Syllogistic Reasoning, Content Effects, and Belief Bias: Progress Report}
\author{Abdullah Saleem \\
	SEECS, NUST \\
	\texttt{i220882@nu.edu.pk}
	\and
	Ahsan Zaidi \\
	SEECS, NUST \\
	\texttt{azaidi.bsai24seecs@seecs.edu.pk}
	\and
	Mehwish Fatima \\
	SEECS, NUST \\
	\texttt{mehwish.fatima@seecs.edu.pk}}
\date{}

\begin{document}
\maketitle

\begin{abstract}
This document summarizes the research progress for Week 1 and Week 2. Week 1 covers background on syllogistic reasoning, content effects, belief bias, and failure modes/mitigation strategies. Week 2 presents advanced dataset exploration and NLP analysis, including classical syllogistic form distributions, parse-tree complexity, and key research questions on linguistic properties, content effects, and believability correlations.
\end{abstract}

\section{Syllogistic Reasoning}

Syllogistic reasoning is a form of deductive reasoning in which a conclusion is drawn from two premises that share a common ``middle'' term. Each premise relates two categories, and the task is to judge whether a given conclusion follows necessarily from the premises---regardless of whether the statements are actually true in the real world.

A classic categorical syllogism has the structure:

Premise 1: All A are B.

Premise 2: All B are C.

Conclusion: All A are C.

Concretely:

All mammals are warm-blooded. \\
All dogs are mammals. \\
Therefore, all dogs are warm-blooded. (Valid)

Validity in syllogistic reasoning depends purely on logical form---the structure of quantifiers (``all,'' ``some,'' ``no'') and the arrangement of terms---not on whether the content is realistic or familiar. This is what makes syllogisms useful for studying human reasoning: researchers can hold logical form constant while varying content, and see whether people's judgments track logic or something else.

This is also why syllogistic reasoning is central to debates about whether humans (and increasingly, LLMs) reason using formal logical rules, or instead rely on heuristics, mental models, or prior beliefs---which is exactly where content effects and belief bias come in.

\section{Content Effect}

A content effect occurs when the believability, familiarity, or real-world meaning of the terms in a syllogism influences how easily or accurately people reason about it, even though the logical structure is identical. In principle, logical validity should be insensitive to content, but in practice, people reason more accurately and more quickly with concrete, familiar content than with abstract or unfamiliar content.

Example---same logical form, different content:

Abstract version:
All A are B. All B are C. Therefore, all A are C.

Concrete version:
All roses are flowers. All flowers need water. Therefore, all roses need water.

Both arguments share the exact same valid structure, but people typically find the concrete, meaningful version easier to evaluate correctly than the abstract A/B/C version. This is the content effect: performance and reasoning strategy shift with subject matter, not just logical form.

Content effects are usually explained by the idea that people reason using mental models or real-world knowledge to fill in meaning, rather than applying abstract formal rules---concrete content gives reasoners more scaffolding to construct and check a mental model of the situation.

\section{Belief Bias}

Belief bias is the tendency to judge a conclusion as logically valid or invalid based on whether it is believable or consistent with prior knowledge, rather than on whether it actually follows logically from the premises. It is a specific and well-studied type of content effect, and it is one of the strongest, most replicated findings in reasoning research.

Belief bias becomes especially visible when logical validity and believability conflict. This produces four combinations:

\begin{itemize}
	\item Conclusion believable / Logically valid: Easy -- most people agree.
	\item Conclusion believable / Logically invalid: Hard -- often wrongly accepted.
	\item Conclusion unbelievable / Logically valid: Hard -- often wrongly rejected.
	\item Conclusion unbelievable / Logically invalid: Easy -- most people agree.
\end{itemize}

The most diagnostic case is valid-but-unbelievable, since it most clearly separates logic from belief:

Premise 1: All flowers need water. \\
Premise 2: Roses need water. \\
Conclusion: Roses are flowers.

And the reverse---a logically valid syllogism with an unbelievable conclusion, which people often incorrectly reject:

Premise 1: All mammals can walk. \\
Premise 2: Whales are mammals. \\
Conclusion: Whales can walk.

Here the conclusion is logically valid given the premises, but because it clashes with real-world knowledge (whales swim, not walk), people frequently---and incorrectly---judge the argument as invalid. This mismatch between logical validity and subjective believability is the hallmark of belief bias.

Belief bias is typically explained by dual-process accounts of reasoning: a fast, intuitive system that evaluates conclusions based on prior belief and plausibility, and a slower, more effortful analytic system capable of genuine logical evaluation. Belief bias emerges when the intuitive system dominates, especially under time pressure, cognitive load, or low motivation to engage in careful analysis.

\section{Why Do LLMs Fail on Logical Tasks?}

When you look at why these models completely mess up on logic, it basically comes down to several issues: semantic entanglement, form heuristics, structural blindspots, and sensitivity to layout.

\paragraph{Semantic Entanglement}
LLMs are trained to predict the next word by reading large text corpora. They learn real-world facts and correlations extremely well, and those learned facts can override pure logical reasoning. Mechanistically, some analyses show that ``factual'' attention heads and hidden directions can leak real-world knowledge into the residual stream, which corrupts pure symbol-tracking required for formal logic.

\paragraph{Atmosphere / Form Heuristics}
On zero-shot setups, models often rely on superficial cues rather than computing symbolic structure. One common shortcut is the Atmosphere Heuristic: models pick up quantifier ``vibes'' like ``Some'' or ``No'' and tend to mirror them in conclusions, producing superficially plausible but logically incorrect outputs.

\paragraph{Nothing-Follows Blindspot}
Many syllogisms correctly yield ``nothing follows'' (no valid conclusion). Pre-trained LLMs have a strong bias to produce some relation rather than admit neutrality, which makes them fail on neutral/invalid test cases.

\paragraph{Figural / Layout Sensitivity}
Models are highly sensitive to token ordering and prompt layout. When variables are presented in a particular sequence ($A \rightarrow B \rightarrow C$), models prefer conclusions matching that order; when the order is changed, their behavior can flip. This shows reliance on sequential text patterns rather than context-free logical form.

\section{Existing Mitigation Approaches}

Researchers have tried several fixes across prompting, fine-tuning, and inference-time interventions:

\begin{itemize}
	\item Translate-and-Explain (Chain-of-Thought): Prompt models to translate natural language into formal symbolic logic before solving. Models are often good at producing formal translations, which can help but do not fully eliminate downstream reasoning errors.
	\item Supervised Fine-Tuning on Pseudo-Words: Replacing nouns with random pseudo-words (e.g., "schmeeft") has been shown to help detach grammatical structure from real-world meaning and reduce belief-driven errors in some settings.
	\item Activation Steering (Inference Interventions): Identify hidden directions correlated with world knowledge and subtract or modulate them during inference to reduce factual leakage and force reliance on logical form.
	\item Conditional Dynamic Steering (K-CAST): Apply a runtime, k-NN-based steering magnitude calibrated to the prompt's degree of factual/logic conflict. This dynamic approach avoids the instability of a one-size-fits-all steering vector.
\end{itemize}

\section{Open Research Challenges}

Even with these approaches, several challenges remain:

\begin{itemize}
	\item The Scaling Limit Paradox: Larger models improve somewhat but plateau, and in some cases become more confidently wrong because they internalize human reasoning errors.
	\item Multi-Step Extrapolation Wall: Models fine-tuned on two-premise syllogisms fail to generalize to longer deductive chains (3--4 premises).
	\item Cross-Lingual Knowledge Mess: Mixing multilingual context during retrieval or prompt augmentation can introduce noisy cross-language signals that hurt logical reasoning.
\end{itemize}

\section{Week 2 Report}

\subsection{Syllogistic Forms Distribution}
Our advanced NLP analysis using spaCy successfully scanned all 960 samples and mapped their quantifier layouts back to their historical, classical Aristotelian logic names. The dataset is surprisingly diverse and doesn't just stick to simple setups.

The single most common structure in the dataset is the Felapton / Fesapo mood, which accounts for 11.56\% of the samples (111 instances). This form uses a no + all $\rightarrow$ some setup, which is a highly complex logical track because it forces a broad universal negative premise to transition into a specific, particular negative conclusion. The classic Barbara structure (all + all $\rightarrow$ all) comes in second at 10.83\% (104 samples), serving as the baseline for clean, straightforward transitivity. The rest of the top distributions are split between Darii (all + some $\rightarrow$ some) at 9.69\%, Ferio / Festino (no + some $\rightarrow$ some) at 9.38\%, and Barbari / Bramantip at 8.96\%.

The fact that the dataset is so evenly spread across these tricky moods---rather than just spamming basic Barbara problems---shows it is intentionally built to test if cognitive models can genuinely follow shifting set boundaries or if they are just gaming surface-level word patterns.

\subsection{Syntactic Analysis (Parse-Tree Complexity)}
To evaluate the grammatical difficulty of the text, we used spaCy to construct linguistic parse-trees for every sentence and measured their average dependency depth. The average parse-tree depth came out to 4.29 levels deep.

In practical terms, a depth of over 4 layers means these sentences are written in a highly natural, human-like cadence with plenty of nested sub-clauses, relative modifiers, and passive phrasing, rather than short, robotic ``A is B'' templates.

The maximum sentence complexity encountered in the dataset reached a staggering 9 levels deep. When a sentence branches out 9 layers into a grammatical tree maze, it creates an incredibly high syntactic load. For a transformer-based Large Language Model, tracking subject-predicate variables across 9 deep structural branches requires intense attention weight coordination, making these specific sentences the ultimate blindspots where text-parsing confusion can easily corrupt the model's logical path.

\subsection{Research Question 1 (RQ1): Which linguistic properties correlate with reasoning difficulty?}
Based on our exploratory analysis, two main linguistic properties stand out as major drivers of reasoning difficulty for LLMs: high parse-tree dependency depth and negation density.

First, when the parse-tree depth spikes up toward the maximum of 9 levels, the text becomes structurally tangled. The model has to carry token representations across multiple layers of sub-clauses before it even reaches the main verb or object. This distance makes it easy for the model's attention mechanism to experience ``recency bias'' or drop key variable relationships entirely.

Second, our negation analysis found 691 negations in the premises and 520 in the conclusions, with nearly 8\% of the dataset containing complex, multi-negation structures (like ``It is not true that any tree is green'' or ``cannot be classified as''). In cognitive science, multiple negations are famous for destroying an agent's mental model. An LLM has to constantly invert its internal vector directions every time it reads a negative token (like not, never, cannot). When these negative tokens are stacked together, the model frequently suffers from operation-fusing errors, failing to compute the correct deduction simply because it got lost in the negative grammar.

\subsection{Research Question 2 (RQ2): Are some syllogistic forms more susceptible to content effects?}
Yes, the data strongly suggests that specific syllogistic forms---specifically Felapton / Fesapo (no + all $\rightarrow$ some) and Ferio / Festino (no + some $\rightarrow$ some)---are significantly more vulnerable to content effects than simpler structures like Barbara.

The reason these forms are so fragile when content changes comes down to their mixed-quantifier design. They force the reasoning engine to jump mid-argument from absolute, sweeping claims (``All'' or ``No'') to localized, partial claims (``Some... not'').

When a structure is already this mentally taxing, introducing a counterfactual content trap (like an Unbelievable Valid scenario where the logic is correct but says something absurd like ``some fish are reptiles'') triggers an immediate error signal in the model's pre-trained semantic weights. Because the structural logic circuit of a Felapton or Ferio form is already complex and fragile, the competing factual activations from the model's internet memory can easily override the transformer's internal middle-term suppression mechanism, causing the model to break down and agree with its biases instead of the math.

\subsection{RQ3: Does believability correlate with validity labels?}

\subsubsection*{Direct Answer (Student-Style Language)}
No, believability does not correlate with validity labels in this dataset. When we ran the cross-reference check between the logical mathematics (validity) and the real-world truth (plausibility), the dataset turned out to be almost perfectly balanced into four equal quadrants of 25\% each (240 samples per quadrant, out of 960 total).

Specifically, 48.75\% of all logically invalid problems are wrapped in highly believable text (the Believable-Invalid trap), while exactly 50.00\% of all logically valid problems are wrapped in completely absurd, counterfactual text (the Unbelievable-Valid test).

\subsubsection*{What this means}
Because the distribution is practically a perfect 50/50 split across the board, the Pearson correlation or general dependency between a statement being factually believable and it being logically valid is essentially 0.

The dataset creators intentionally engineered it this way so that a Large Language Model cannot use real-world plausibility as a ``cheat code'' or statistical shortcut to guess the right answer. The data completely strips away any correlation, forcing the model to ignore its pre-trained semantic world-knowledge and rely strictly on structural, context-free logic to get the classification right.

\bibliographystyle{acl_natbib}
\bibliography{anthology,custom}

\end{document}
